\section*{\centering{Abstract}}
\addcontentsline{toc}{section}{Abstract}
\vspace{0.5cm}
In most colleges the class timetable is published as a multi-page PDF that every student and teacher
has to read by hand. Finding the next class, a free room or a teacher's location means scanning dense
grids full of abbreviations, lab sub-groups and footnote legends, and every revision of the timetable
has to be re-checked manually. Clashes between teachers or rooms are usually discovered only when they
happen.

This mini project builds \textbf{Dynamic Timetable}, a role-based timetable monitoring system made of
a cross-platform mobile and web app (React Native with Expo) and a central back end (Python FastAPI
and PostgreSQL). The Admin uploads the official timetable as PDF, Excel, Word or a photo. A
background worker extracts every class cell with geometry-aware table parsing and OCR, expands the
per-page legends, separates lab sub-groups, validates times and flags anything uncertain for review.
Students, professors, the HOD and the Principal then see a dashboard (current class, next class,
free periods) and can ask questions in plain English such as ``Do I have classes after lunch on
Tuesday?'', ``How many labs do I have this week?'' or ``Which day is my busiest?''. A deterministic
compositional planner turns each question into a typed query plan (an operation such as list, count,
first/last, total hours or day-wise summary, combined with filters, a date or week, a time window and
an ordering), which the server compiles into parameterised queries with the user's role scope added.
It never generates SQL from text and cannot be talked into reading other users' data. Institutional and personal timetables are
kept strictly separate, and the system automatically detects professor, room and batch clashes.

On the college's real 10-page timetable PDF the system extracted 453 class entries, of which 452
were verified automatically; the one remaining entry has no time printed in the PDF. It also reported
exactly the two genuine clashes present in the document. The work is supported by 174 automated
back-end tests, 20 mobile unit tests and end-to-end browser tests (24/24 and 21/21 passing).
Measured single-user search latency was about 37--47~ms (median) on a 4-core development machine.
The system has been deployed as a free online demo. Its limitations are that it has not been load
tested on production hardware, push notifications need a store build, and extraction accuracy has
not yet been measured against a hand-labelled ground truth.

\vspace{0.5cm}
\noindent\textbf{Keywords:} timetable extraction, table recognition, OCR, natural-language interface,
intent classification, role-based access control, conflict detection, React Native, FastAPI,
PostgreSQL.
\newpage
